\documentclass[../book.tex]{subfiles}
\begin{document}

    \section{微分方程}

    \subsection{微分方程的概念}

    \subsubsection{微分方程及其阶}

    含有未知函数、未知函数的导数（或微分）与自变量之间关系的方程，称为\textbf{微分方程}，一般写成
    \[
        F[x, y, y',\cdots, y^{(n)}] = 0 \quad \text{或} \quad y^{(n)} = f[x, y, y',\cdots, y^{(n-1)}]
    \]

    \textbf{微分方程的阶}：微分方程中出现的未知函数的最高阶导数的阶数，称为该微分方程的\textbf{阶}

    \subsubsection{常微分方程}

    未知函数是一元函数（即只有一个自变量）的微分方程，称为\textbf{常微分方程}（Ordinary Differential Equation, ODE）。如$y''' - y'' + 6y' = 0,\quad y\mathrm{d}x - (x - \sqrt{x^2 + y^2})\mathrm{d}y = 0$

    \subsubsection{线性微分方程}

    如果微分方程中，未知函数 $y$ 及其各阶导数 $y', y'', \ldots, y^{(n)}$ 都是一次的（即它们的幂次都是 $1$），且不含这些函数的乘积或复合，则称为\textbf{线性微分方程}。

    \begin{enumerate}
        \item \textbf{$n$ 阶线性微分方程}：

        一般形式：
        \[
            a_n(x)y^{(n)} + a_{n-1}(x)y^{(n-1)} + \cdots + a_1(x)y' + a_0(x)y = f(x)
        \]

        其中系数 $a_n(x), a_{n-1}(x), \ldots, a_0(x)$ 和 $f(x)$ 都是 $x$ 的已知函数，且 $a_n(x) \neq 0$。

        \item \textbf{$n$ 阶常系数线性微分方程}：

        若上述方程中所有系数 $a_n, a_{n-1}, \ldots, a_0$ 都是常数（与 $x$ 无关），则称为\textbf{常系数}线性微分方程。一般形式：
        \[
            a_ny^{(n)} + a_{n-1}y^{(n-1)} + \cdots + a_1y' + a_0y = f(x)
        \]

        通常可化为 $a_n = 1$ 的标准形式：
        \[
            y^{(n)} + p_{n-1}y^{(n-1)} + \cdots + p_1y' + p_0y = g(x)
        \]

        \item \textbf{$n$ 阶齐次线性微分方程}：

        若 $f(x) \equiv 0$（即右端自由项为零），则称为\textbf{齐次}线性微分方程。一般形式：
        \[
            a_n(x)y^{(n)} + a_{n-1}(x)y^{(n-1)} + \cdots + a_1(x)y' + a_0(x)y = 0
        \]

        \item \textbf{$n$ 阶非齐次线性微分方程}：

        若 $f(x) \not\equiv 0$（即右端自由项不为零），则称为\textbf{非齐次}线性微分方程。一般形式：
        \[
            a_n(x)y^{(n)} + a_{n-1}(x)y^{(n-1)} + \cdots + a_1(x)y' + a_0(x)y = f(x)
        \]
    \end{enumerate}

    \begin{analysisbox}[分类举例]
        \begin{itemize}
            \item $y'' + 3y' + 2y = 0$：二阶常系数齐次线性微分方程
            \item $y'' + 3y' + 2y = e^x$：二阶常系数非齐次线性微分方程
            \item $y'' + xy' + y = 0$：二阶变系数齐次线性微分方程
            \item $y'' + xy' + y = \sin x$：二阶变系数非齐次线性微分方程
        \end{itemize}
    \end{analysisbox}

    \subsubsection{微分方程的解}

    代入微分方程能使方程成为恒等式的函数，称为该微分方程的\textbf{解}。微分方程解的图形称为\textbf{积分曲线}

    \subsubsection{微分方程的通解}

    如果微分方程的解中含有的\textbf{独立常数的个数等于微分方程的阶数}，则称此解为该微分方程的\textbf{通解}。

    \subsubsection{初始条件与特解}

    \begin{enumerate}
        \item \textbf{初始条件}：用来确定通解中任意常数的附加条件
        \item \textbf{特解}：由初始条件确定了通解中任意常数后得到的解
    \end{enumerate}

    \subsection{一阶微分方程的求解}

    \subsubsection{可分离变量型微分方程}

    \begin{enumerate}
        \item \textbf{直接可分离}
        \begin{itemize}
            \item \textbf{标准形式}：$\dfrac{\mathrm{d}y}{\mathrm{d}x} = f(x)g(y)$
            \item \textbf{特征}：方程可分解为只含 $x$ 的函数与只含 $y$ 的函数的乘积。
            \item \textbf{解法}：分离变量后两边积分
            \[
                \int \frac{1}{g(y)}\,\mathrm{d}y = \int f(x)\,\mathrm{d}x + C
            \]
        \end{itemize}


        \item \textbf{换元后可分离}
        \begin{itemize}
            \item \textbf{标准形式}：$\dfrac{\mathrm{d}y}{\mathrm{d}x} = f(ax + by + c) \quad (a, b \text{ 不全为零})$
            \item \textbf{特征}：方程右端是 $ax + by + c$ 的函数。
            \item \textbf{解法}：令 $u = ax + by + c$，则
            \[
                \frac{\mathrm{d}u}{\mathrm{d}x} = a + b\frac{\mathrm{d}y}{\mathrm{d}x} = a + bf(u)
            \]
            分离变量得 $\displaystyle\frac{\mathrm{d}u}{a + bf(u)} = \mathrm{d}x$，积分后回代即可。
        \end{itemize}
    \end{enumerate}

    \subsubsection{齐次型微分方程}

    \begin{itemize}
        \item \textbf{标准形式}：$\displaystyle\frac{\mathrm{d}y}{\mathrm{d}x} = \varphi\left(\frac{y}{x}\right)$
        \item \textbf{特征}：方程中 $x$ 和 $y$ 的地位对称，右端是 $\dfrac{y}{x}$ 的函数
        \item \textbf{解法}：令 $u = \dfrac{y}{x}$（即 $y = xu$），则 $\dfrac{\mathrm{d}y}{\mathrm{d}x} = u + {\color{red}{x\dfrac{\mathrm{d}u}{\mathrm{d}x}}}$，代入后得到关于 $u$ 的可分离变量方程：
        \[
            x\frac{\mathrm{d}u}{\mathrm{d}x} = \varphi(u) - u \implies \frac{\mathrm{d}u}{\varphi(u)-u} = \frac{\mathrm{d}x}{x}
        \]
    \end{itemize}

    \subsubsection{一阶线性微分方程}

    \begin{itemize}
        \item \textbf{标准形式}：$y' + p(x)y = q(x)$ 或 $x' + p(y)x = q(y)$
        \item \textbf{特征}：$y$ 和 $y'$ 都是一次的（线性），系数 $p(x)$、$q(x)$ 仅与 $x$ 有关
        \item \textbf{通解公式}：
        \[
            y = e^{-\int p(x)\,\mathrm{d}x}\left[\int q(x)e^{\int p(x)\,\mathrm{d}x}\,\mathrm{d}x + C\right]
        \]
        \item \textbf{定积分形式}：
        \[
            y(x) = e^{-\int_{x_0}^{x} p(t)\,\mathrm{d}t}\left[\int_{x_0}^{x} q(t)e^{\int_{x_0}^{t} p(s)\,\mathrm{d}s}\,\mathrm{d}t + C\right]
        \]
    \end{itemize}

    \begin{analysisbox}[推导：一阶线性微分方程通解（积分因子法）]
        \textbf{分析}：构造积分因子 $\mu(x)$，把左端整理成一个乘积的导数，从而可直接积分。

        \textbf{解}：设积分因子
        \[
            \mu(x)=e^{\int p(x)\,\mathrm{d}x}
        \]
        则 $\mu'(x)=p(x)\mu(x)$。原方程
        \[
            y' + p(x)y = q(x)
        \]
        两边同乘 $\mu(x)$，得
        \[
            \mu(x)y' + \mu(x)p(x)y = \mu(x)q(x)
        \]
        由 $\mu'=p\mu$，左端可写为
        \[
            (\mu y)'=\mu q
        \]

        两边积分
        \[
            \mu y=\int \mu q\,\mathrm{d}x + C
        \]
        解出 $y$：
        \[
            y=e^{-\int p(x)\,\mathrm{d}x}\left[\int q(x)e^{\int p(x)\,\mathrm{d}x}\,\mathrm{d}x+C\right]
        \]
    \end{analysisbox}

    \subsubsection{伯努利方程}

    \begin{itemize}
        \item \textbf{标准形式}：$\displaystyle\frac{\mathrm{d}y}{\mathrm{d}x} + p(x)y = q(x)y^n$ 或 $\displaystyle\frac{\mathrm{d}x}{\mathrm{d}y} + p(y)x = q(y)x^n \quad (n \neq 0, 1)$
        \item \textbf{解法}：
        \begin{circlenum}
            \item 恒等变形：方程两边除以 $y^n$，得 $y^{-n}\dfrac{\mathrm{d}y}{\mathrm{d}x} + p(x)y^{1-n} = q(x)$
            \item 令 $z = y^{1-n}$，得 $\dfrac{\mathrm{d}z}{\mathrm{d}x} = (1-n)y^{-n}\dfrac{\mathrm{d}y}{\mathrm{d}x}$，则 $\dfrac{1}{1-n}\dfrac{\mathrm{d}z}{\mathrm{d}x} + p(x)z = q(x)$
            \item 代入得{\color{red}{一阶线性方程}}：$\dfrac{\mathrm{d}z}{\mathrm{d}x} + (1-n)p(x)z = (1-n)q(x)$
        \end{circlenum}
    \end{itemize}

    \subsubsection{二阶可降阶微分方程（用换元法化为一阶方程）}

    \begin{enumerate}
        \item $y'' = f(x, y')$型
        \begin{itemize}
            \item \textbf{特征}：方程中不显含 $y$
            \item \textbf{解法（{\color{red}{赶尽杀绝$y$}}）}
            \begin{circlenum}
                \item 令 $p = y' = p(x)$，则 $y'' = p' = \dfrac{dp}{dx}$，则原方程变为一阶方程 $\displaystyle\frac{dp}{dx} = f(x, p)$
                \item 若求得其通解为 $p = \varphi(x, C_1)$，即 $y' = \varphi(x, C_1)$，则原方程的通解为 $\displaystyle y = \int \varphi(x, C_1)\mathrm{d}x + C_2$
            \end{circlenum}
        \end{itemize}
        \item $y'' = f(y, y')$型
        \begin{itemize}
            \item \textbf{特征}：方程中不显含自变量 $x$
            \item \textbf{解法（{\color{red}{斩草除根$x$}}）}
            \begin{circlenum}
                \item 令 $p = y' = p(y)$，则 ${\color{red}{y'' = \dfrac{dp}{dx}}} = \dfrac{dp}{dy} \cdot \dfrac{dy}{dx} = p\dfrac{dp}{dy}$，则原方程变为一阶方程 $\displaystyle p\frac{dp}{dy} = f(y, p)$
                \item 若求得其通解为 $p = \varphi(y, C_1)$，即 $y' = \varphi(y, C_1)$，则{\color{red}{分离变量积分}}得 $\dfrac{dy}{\varphi(y, C_1)} = dx$
                \item 两边积分得 $\displaystyle\int \dfrac{dy}{\varphi(y, C_1)} = x + C_2$，即可求得原方程的通解
            \end{circlenum}
        \end{itemize}
        \item $y'' = f(y')$型
        \begin{itemize}
            \item \textbf{特征}：方程中不显含 $y$
            \item \textbf{解法（{\color{red}{赶尽杀绝$y$}}）}
            \begin{circlenum}
                \item 令 $p = y'$，则 $y'' = p' = \dfrac{dp}{dx}$，原方程变为 $\dfrac{dp}{dx} = f(p)$
                \item 分离变量：$\dfrac{dp}{f(p)} = dx$，积分得 $p = \varphi(x, C_1)$，即 $y' = \varphi(x, C_1)$
                \item 再积分得原方程通解：$\displaystyle y = \int \varphi(x, C_1)\mathrm{d}x + C_2$
            \end{circlenum}
        \end{itemize}
    \end{enumerate}

    \subsection{高阶线性微分方程的求解}

    \subsubsection{二阶常系数齐次线性微分方程}

    \begin{enumerate}
        \item \textbf{标准形式}：$y'' + py' + qy = 0$（$p, q$ 为常数）
        \item \textbf{解的结构}：若 $y_1(x), y_2(x)$ 是方程 $y'' + py' + qy = 0$ 的两个解，且 $\dfrac{y_1}{y_2} \neq$ 常数，则称它们为线性无关解，且 $y = C_1y_1 + C_2y_2$ 是方程的通解
        \item $\color{red}{\bigstar}$\textbf{通解}（特征方程法）：
        \begin{circlenum}
            \item 写出特征方程：$r^2 + pr + q = 0$
            \item 根据特征根的情况写出通解：
            \begin{itemize}
                \item $p^2 - 4q > 0$，设 $r_1, r_2$ 是特征方程的\textbf{两个不等实根}，即 $r_1 \neq r_2$，可得其通解为：
                \[
                    y = C_1e^{r_1x} + C_2e^{r_2x}
                \]
                \item $p^2 - 4q = 0$，设 $r$ 是特征方程的\textbf{二重实根}，即 $r_1 = r_2 = r$，可得其通解为：
                \[
                    y = (C_1 + C_2x)e^{rx}
                \]
                \item $p^2 - 4q < 0$，设 $\alpha \pm i\beta$ 是特征方程的\textbf{一对共轭复根}，可得其通解为：
                \[
                    y = e^{\alpha x}(C_1\cos\beta x + C_2\sin\beta x)
                \]
            \end{itemize}
        \end{circlenum}
    \end{enumerate}

    \subsubsection{二阶常系数非齐次线性微分方程}

    \begin{enumerate}
        \item \textbf{标准形式}：$y'' + py' + qy = f(x) \quad (p, q \text{ 为常数}, f(x) \neq 0)$，$f(x)$ 称为自由项
        \item \textbf{解的结构}
        \begin{itemize}
            \item \textbf{叠加原理}：若 $f(x) = f_1(x) + f_2(x)$，且 $y_1^*$ 是 $y'' + py' + qy = f_1(x)$ 的解，$y_2^*$ 是 $y'' + py' + qy = f_2(x)$ 的解，则 $y_1^* + y_2^*$ 是原方程的解
            \item 若 $y_1^*, y_2^*$ 是同一非齐次方程的两个特解，则 $y_1^* - y_2^*$ 是对应齐次方程的解
        \end{itemize}

        \item \textbf{通解}
        \[
            y = Y + y^*
        \]
        其中 $Y$ 是对应齐次方程的通解，$y^*$ 是非齐次方程的特解。

        \item \textbf{特解的设定}
        \begin{itemize}
            \item \textbf{待定系数法（{\color{red}{一看二算三比较}}）}

            设 $P_n(x), P_m(x)$ 为 $x$ 的 $n$ 次多项式

            \begin{enumerate}
                \item 当自由项 $f(x) = P_n(x)e^{\alpha x}$，即$y'' + py' + qy = P_n(x)e^{\alpha x}$，特解设为
                \[
                    y^* = e^{\alpha x} Q_n(x)x^k
                \]
                其中
                \[
                    \begin{cases}
                        e^{\alpha x} \text{照抄} \\[6pt]
                        Q_n(x) = a_nx^n + a_{n-1}x^{n-1} + \cdots + a_1x + a_0 \text{与 } P_n(x) \text{ 同次多项式，系数待定} \\
                        \emph{e.g.} \text{ 若} P_1(x) = 2x, \text{ 则} Q_1(x) = ax + b  \\[6pt]
                        k = \begin{cases}
                                0, & \alpha \text{ 不是特征根，} \alpha \neq r_{1,2} \\
                                1, & \alpha \text{ 是单特征根，} \alpha = r_1 \text{或 } \alpha = r_{2}(r_1 \neq r_2)   \\
                                2, & \alpha \text{ 是二重特征根，} \alpha = r_1 = r_2
                        \end{cases}
                    \end{cases}
                \]

                \item 当自由项 $f(x) = e^{\alpha x}[P_m(x)\cos\beta x + P_n(x)\sin\beta x]$，即$y'' + py' + qy = e^{\alpha x}[P_m(x)\cos\beta x + P_n(x)\sin\beta x]$，特解设为
                \[
                    y^* = e^{\alpha x}[Q_l^{(1)}(x)\cos\beta x + Q_l^{(2)}(x)\sin\beta x]x^k
                \]
                其中
                \[
                    \begin{cases}
                        e^{\alpha x} \text{照抄} \\[6pt]
                        l = \max\{m, n\} Q_l^{(1)}(x), Q_l^{(2)}(x) \text{ 都是 } m \text{ 次多项式，系数待定} \\[6pt]
                        k = \begin{cases}
                                0, &  \alpha + \beta i \neq r_{1, 2} \\
                                1, &  \alpha + \beta i = r_{1, 2}
                        \end{cases}
                    \end{cases}
                \]
            \end{enumerate}

            \item \textbf{微分算子法}
            \begin{enumerate}
                \item \textbf{基本思想}：引入微分算子 $D = \dfrac{d}{dx}$，将微分方程转化为代数运算。
                \item \textbf{算子定义}：
                \begin{enumerate}
                    \item $D = \displaystyle\frac{d}{dx}, Dy = y', D^2 = \displaystyle\frac{d^2}{dx^2}, D^2y = y''$，$D^ny = y^{(n)}$
                    \item 方程 $y'' + py' + qy = f(x)$ 可写成 $(D^2 + pD + q)y = f(x)$
                    \item 记 $F(D) = D^2 + pD + q$，则 $F(D)y = f(x)$，此时它的一个特解为 $y^* = \dfrac{1}{F(D)}f(x)$
                    \item $D$表示求导的条件下，$\displaystyle\frac{1}{D}$表示积分。如$D\sin x = \cos x, \dfrac{1}{D}\sin x = -\cos x$
                \end{enumerate}
                \item \textbf{算子法常见类型}：

                \begin{enumerate}
                    \item $\dfrac{1}{F(D)}e^{\alpha x} $型

                    \begin{circlenum}
                        \item 若 $F(\alpha) \neq 0$，有 $y^* = \dfrac{e^{\alpha x}}{F(\alpha)}$
                        \item 若 $F(\alpha) = 0$，$\alpha$ 是 $k$ 重特征根，有 $y^* = \dfrac{x^ke^{\alpha x}}{F^{(k)}(\alpha)}$
                    \end{circlenum}

                    \item $\dfrac{1}{D^2 + q}\cos \beta x$ 或 $\dfrac{1}{D^2 + q}\sin \beta x$ 型

                    \begin{circlenum}
                        \item 若 $(D^2 + q)\Big |_{D = \beta i} \neq 0$，有
                        \[
                            y^* = \dfrac{1}{q-\beta^2}\cos\beta x \quad \text{或} \quad y^* = \dfrac{1}{q-\beta^2}\sin\beta x
                        \]
                        \item 若 $(D^2 + q)\Big |_{D = \beta i} = 0$（即 $q = \beta^2$），有
                        \[
                            y^* = x\cdot\dfrac{1}{(D^2 + q)'}\cos\beta x \quad \text{或} \quad y^* = x\cdot\dfrac{1}{(D^2 + q)'}\sin\beta x
                        \]
                    \end{circlenum}

                    \item $\dfrac{1}{F(D)}\cos \beta x$ 或 $\dfrac{1}{F(D)}\sin \beta x$ 型

                    若 $F(D) = D^2 + pD + q$，则取 $F(D)\Big |_{D^2 = (\beta i)^2} = pD + q - \beta^2$
                    \begin{circlenum}
                        \item
                        \[
                            y^* = \dfrac{1}{F(D)}\cos \beta x = \dfrac{(q-\beta^2)\cos \beta x + p\beta\sin \beta x}{(q-\beta^2)^2 + p^2\beta^2}
                        \]
                        \item
                        \[
                            y^* = \dfrac{1}{F(D)}\sin \beta x = \dfrac{(q-\beta^2)\sin \beta x - p\beta\cos \beta x}{(q-\beta^2)^2 + p^2\beta^2}
                        \]
                    \end{circlenum}

                    \item $\dfrac{1}{F(D)}(x^k + a_1x^{k-1}+\cdots+a_{k-1}x + a_k)$ 型

                    将 $\dfrac{1}{F(D)}$ 按 $D$ 的升幂级数展开：
                    \[
                        y^* = \dfrac{1}{F(D)}(x^k + a_1x^{k-1}+\cdots+a_{k-1}x + a_k) = Q_k(D)(x^k + a_1x^{k-1}+\cdots+a_{k-1}x + a_k)
                    \]
                    其中 $Q_k(D)$ 是将 $\dfrac{1}{F(D)}$ 展开并截断到 $D^k$ 项（常借助$\dfrac{1}{1-x}$的$k$次泰勒多项式$1 + x + x^2 + \cdots + x^k$），即 $b_0 + b_1D + b_2D^2 + \cdots + b_kD^k$
                    \begin{analysisbox}[Example: 多项式型自由项]
                        已知 $y'' + y' = x^2 + 1$，求$y^*$

                        \textbf{解}：将原方程写成算子形式：$F(D)y = (D^2 + D)y = x^2 + 1$

                        特解为：
                        \[
                            y^* = \dfrac{1}{D^2 + D}(x^2 + 1)
                        \]

                        将 $\dfrac{1}{D^2 + D}$ 按 $D$ 的升幂级数展开：
                        \[
                            \frac{1}{D^2 + D} = \frac{1}{D}\cdot\frac{1}{1 + D} = \frac{1}{D}(1 - D + D^2 - D^3 + \cdots)
                        \]

                        截断到 $D^2$ 项（因为右端是二次多项式，更高次项作用后为零）：
                        \[
                            Q_2(D) = \frac{1}{D} - 1 + D
                        \]

                        因此：
                        \[
                            \begin{aligned}
                                y^*
                                &= \left(\frac{1}{D} - 1 + D\right)\left(x^2+1\right) \\
                                &= \frac{1}{D}(x^2 + 1) - \left(x^2+1\right) + D\left(x^2+1\right) \\
                                &= \frac{x^3}{3} + x - x^2 - 1 + 2x \\
                                &= \frac{x^3}{3} - x^2 + 3x - 1
                            \end{aligned}
                        \]
                    \end{analysisbox}

                    \item $\dfrac{1}{F(D)}e^{\alpha x}v(x)$ 型（移位定理）

                    \textbf{移位定理}：$\dfrac{1}{F(D)}[e^{\alpha x}v(x)] = e^{\alpha x}\dfrac{1}{F(D+\alpha)}v(x)$

                    \[
                        y^* = e^{\alpha x}\dfrac{1}{F(D+\alpha)}v(x)
                    \]
                    将问题转化为 $\dfrac{1}{F(D+\alpha)}v(x)$，再用前述方法求解。
                \end{enumerate}
            \end{enumerate}

        \end{itemize}
    \end{enumerate}

    \subsubsection{n（n>2）阶常系数齐次线性微分方程}

    \begin{enumerate}
        \item \textbf{标准形式}：$y^{(n)} + a_1y^{(n-1)} + \cdots + a_{n-1}y' + a_ny = 0$（$a_1, \ldots, a_n$ 为常数）
        \item \textbf{通解}（特征方程法）：
        \begin{itemize}
            \item 写出特征方程：$r^n + a_1r^{n-1} + \cdots + a_{n-1}r + a_n = 0$
            \item 根据特征根的情况写出通解：
            \begin{circlenum}
                \item 若 $r$ 为单实根：对应项 $Ce^{rx}$
                \item 若 $r$ 为 $k$ 重实根：对应项 $(C_1 + C_2x + \cdots + C_kx^{k-1})e^{rx}$
                \item 若 $r$ 为单复根 $\alpha \pm \beta i$：对应项 $e^{\alpha x}(C_1\cos\beta x + C_2\sin\beta x)$
                \item 若 $r$ 为二重复根 $\alpha \pm \beta i$：对应项 $e^{\alpha x}[(C_1 + C_2x)\cos\beta x + (C_3 + C_4x)\sin\beta x]$
            \end{circlenum}
        \end{itemize}
        \item 反求方程的理论基础
        \begin{circlenum}
            \item 如果解中含特解 $\bm{e^{rx}}$，则 $r$ 至少为单实根；如二重根 $(C_1 + C_2x)e^{rx}$，令 $C_1 = 1$，$C_2 = 0$
            \item 如果解中含特解 $\bm{x^{k-1}e^{rx}}$，则 $r$ 至少为 $k$ 重实根
            \item 如果解中含特解 $\bm{e^{\alpha x}\cos\beta x}$ 或 $\bm{e^{\alpha x}\sin\beta x}$，则 $\alpha \pm \beta i$ 至少为单复根
            \item 如果解中含特解 $\bm{e^{\alpha x}x\cos\beta x}$ 或 $\bm{e^{\alpha x}x\sin\beta x}$，则 $\alpha \pm \beta i$ 至少为二重复根
        \end{circlenum}
    \end{enumerate}

    \subsubsection{欧拉方程}

    \begin{itemize}
        \item \textbf{标准形式}：$x^2y'' + pxy' + qy = f(x)$
        \item \textbf{特征}：各项中 $x$ 的幂次与 $y$ 的导数阶数相同
        \item \textbf{解法}
        \begin{circlenum}
            \item 当 $x > 0$ 时，令 $x = e^t$，则 $t = \ln x$，$\dfrac{dt}{dx} = \dfrac{1}{x}$，于是
            \[
                \frac{dy}{dx} = \frac{dy}{dt} \cdot \frac{dt}{dx} = \frac{1}{x}\frac{dy}{dt}
            \]
            \[
                \frac{d^2y}{dx^2} = \frac{d}{dx}\left(\frac{1}{x}\frac{dy}{dt}\right) = -\frac{1}{x^2}\frac{dy}{dt} + \frac{1}{x}{\color{red}{\frac{d^2y}{dt^2}\frac{dt}{dx}}} = \frac{1}{x^2}\left(\frac{d^2y}{dt^2} - \frac{dy}{dt}\right)
            \]
            方程化为关于 $t$ 的常系数线性方程
            \item 当 $x < 0$ 时，令 $x = -e^t$，则 $t = \ln(-x)$，变换公式相同
        \end{circlenum}
    \end{itemize}

    \subsection{微分方程的几何应用}

    \textbf{核心思路}：利用导数的几何意义建立微分方程，常见类型包括：

    \begin{enumerate}
        \item \textbf{追踪问题}
        \begin{itemize}
            \item \textbf{基本特征}：追踪者的速度方向始终指向被追踪者（目标）。
            \item \textbf{追逐曲线}：追踪者运动的轨迹称为追逐曲线

            \item \textbf{建模方法}：
            \begin{circlenum}
                \item 设追踪者在点 $(x, y)$，被追踪者在点 $(X, Y)$
                \item 追踪者的速度方向指向目标：$\dfrac{dy}{dx} = \dfrac{Y - y}{X - x}$
                \item 结合速度关系建立微分方程
            \end{circlenum}
        \end{itemize}
        \item 解题步骤
        \begin{circlenum}
            \item 设所求曲线为 $y = y(x)$
            \item 根据几何条件，利用导数的几何意义建立微分方程
            \item 解微分方程，求出 $y = y(x, C)$
            \item 根据初始条件或边界条件确定常数 $C$
        \end{circlenum}
    \end{enumerate}

    \subsection{微分方程的物理应用}

    \textbf{核心思路}：根据物理定律建立微分方程

    \subsection{基础概念}

    \begin{enumerate}
    \end{enumerate}

    \subsection{结论}

    \begin{enumerate}
    \end{enumerate}

    \subsection{定理}

    \begin{enumerate}
    \end{enumerate}

    \subsection{运算}

    \begin{enumerate}
    \end{enumerate}

    \subsection{公式}

    \begin{enumerate}
    \end{enumerate}

    \subsection{方法总结}

    \begin{enumerate}
    \end{enumerate}

    \subsection{条件转换思路}

    \begin{enumerate}
    \end{enumerate}

    \subsection{理解}

    \begin{enumerate}
    \end{enumerate}

\end{document}
